%%%%%%%%%%%%%%%%%
% This is an example CV created using altacv.cls (v1.6.4, 13 Nov 2021) written by
% LianTze Lim (liantze@gmail.com), based on the
% Cv created by BusinessInsider at http://www.businessinsider.my/a-sample-resume-for-marissa-mayer-2016-7/?r=US&IR=T
%
%% It may be distributed and/or modified under the
%% conditions of the LaTeX Project Public License, either version 1.3
%% of this license or (at your option) any later version.
%% The latest version of this license is in
%%    http://www.latex-project.org/lppl.txt
%% and version 1.3 or later is part of all distributions of LaTeX
%% version 2003/12/01 or later.
%%%%%%%%%%%%%%%%

%% Use the "normalphoto" option if you want a normal photo instead of cropped to a circle
% \documentclass[10pt,a4paper,normalphoto]{altacv}

\documentclass[10pt,a4paper,ragged2e,withhyper]{altacv}

%% AltaCV uses the fontawesome5 package.
%% See http://texdoc.net/pkg/fontawesome5 for full list of symbols.

% Change the page layout if you need to
\geometry{left=1.25cm,right=1.25cm,top=1.5cm,bottom=1.5cm,columnsep=1.2cm}

% The paracol package lets you typeset columns of text in parallel
\usepackage{paracol}


% Change the font if you want to, depending on whether
% you're using pdflatex or xelatex/lualatex
\ifxetexorluatex
  % If using xelatex or lualatex:
  \setmainfont{Lato}
\else
  % If using pdflatex:
  \usepackage[default]{lato}
\fi

% Change the colours if you want to
\definecolor{VividPurple}{HTML}{3E0097}
\definecolor{SlateGrey}{HTML}{2E2E2E}
\definecolor{LightGrey}{HTML}{666666}
% \colorlet{name}{black}
% \colorlet{tagline}{PastelRed}
\colorlet{heading}{VividPurple}
\colorlet{headingrule}{VividPurple}
% \colorlet{subheading}{PastelRed}
\colorlet{accent}{VividPurple}
\colorlet{emphasis}{SlateGrey}
\colorlet{body}{LightGrey}

% Change some fonts, if necessary
% \renewcommand{\namefont}{\Huge\rmfamily\bfseries}
% \renewcommand{\personalinfofont}{\footnotesize}
% \renewcommand{\cvsectionfont}{\LARGE\rmfamily\bfseries}
% \renewcommand{\cvsubsectionfont}{\large\bfseries}

% Change the bullets for itemize and rating marker
% for \cvskill if you want to
\renewcommand{\itemmarker}{{\small\textbullet}}
\renewcommand{\ratingmarker}{\faCircle}

%% Use (and optionally edit if necessary) this .tex if you
%% want to use an author-year reference style like APA(6)
%% for your publication list
\input{pubs-authoryear}

%% Use (and optionally edit if necessary) this .tex if you
%% want an originally numerical reference style like IEEE
%% for your publication list
% \input{pubs-num}

%% sample.bib contains your publications
\addbibresource{sample.bib}

\begin{document}
\name{Álvaro Silva}
\tagline{Ph.D. Candidate in Software Engineering at FEUP, affiliated with Carnegie Mellon University}
% Cropped to square from https://en.wikipedia.org/wiki/Marissa_Mayer#/media/File:Marissa_Mayer_May_2014_(cropped).jpg, CC-BY 2.0
%% You can add multiple photos on the left or right
\photoR{2.5cm}{alvaro_photo_better.jpg}
% \photoL{2cm}{Yacht_High,Suitcase_High}
%\personalinfo{%
  % Not all of these are required!
  % You can add your own with \printinfo{symbol}{detail}
%  \email{amfpsilva@hotmail.com}
%    \orcid{0009-0005-2941-9942}
%  \linkedin{AlvaroFestasSilva}
%       \NewInfoField{gitlab}{\faGitlab}[https://gitlab.com/]
%  \gitlab{ricostynha}
%  \github{ricostynha1}   
%}


\definecolor{linkcolor}{RGB}{0,100,200}

\newcommand{\socialicon}[3]{%
  \href{#1\sociallinkcolor #3}{\textcolor{VividPurple}{#2}}%
}

\newcommand{\sociallinkcolor}{}
\personalinfo{%
  \email{amfpsilva@hotmail.com}
  \orcid{0009-0005-2941-9942}
  \homepage{ricostynha1.github.io/alvaro_silva}
  %\socialicon{https://orcid.org/}{{\faOrcid}}{0009-0005-2941-9942}\quad
  \socialicon{https://www.linkedin.com/in/}{{\faLinkedinIn}}{AlvaroFestasSilva}\quad
  \socialicon{https://gitlab.com/}{{\faGitlab}}{ricostynha}\quad
  \socialicon{https://github.com/}{{\faGithub}}{ricostynha1}
}

\makecvheader

%% Depending on your tastes, you may want to make fonts of itemize environments slightly smaller
\AtBeginEnvironment{itemize}{\small}

%% Set the left/right column width ratio to 6:4.
\columnratio{0.55}

% Start a 2-column paracol. Both the left and right columns will automatically
% break across pages if things get too long.
\begin{paracol}{2}

\cvsection{Experience}
\cvevent{Ph.D Candidate}{FEUP affiliated CMU}{September 2024 -- Ongoing}{Porto, Portugal}
\begin{itemize}
\item Grant CMU Portugal Affiliated Ph.D. Program
\item Research increase usability of formal tools with focus on Dafny
\item Member of the project Verifixer
\end{itemize}
\divider

\cvevent{Guest Assistant}{FEUP}{ Sep 2024 - Ongoing}{Porto, Portugal}
\begin{itemize}
\item Lectured classes of Functional and Logic Programming, Fundaments of Programming in Python and Computing Systems
\end{itemize}
\divider

\cvevent{ASIC Digital Design Staff Enginner}{Synopsys}{Aug 2021 - Sep 2024}{Porto, Portugal}
\begin{itemize}
\item Carry out functional verification for digital IP with focus on Automotive products.
\item Certificated in ISO 26262 for automotive safety products
\item Explore and apply formal verification (using VCFormal).
\item Create assertions to verify properties of the design from specifications.
\end{itemize}
\divider

\cvevent{Researcher}{INESC TEC}{ Oct 2020 - Aug 2021}{Porto, Portugal}
\begin{itemize}
\item Integrated the European investigation project TRUST AI.
\item Granted two research scholarships during that time.    
    
\end{itemize}
\divider
% \divider

% \cvevent{Product Engineer}{Google}{23 June 1999 -- 2001}{Palo Alto, CA}

% \begin{itemize}
% \item Joined the company as employe \#20 and female employee \#1
% \item Developed targeted advertisement in order to use user's search queries and show them related ads
% \end{itemize}

\cvsection{Education}

\cvevent{Integrated Masters in Electrical and Computer Engineering}{Faculty of Engineering of the University of Porto (FEUP)}{Sept 2016 -- June 2021}{ Graduated with \textbf{19}/20}
Minor in Telecommunications 
Specialization in Micro-electronics

\divider

\cvevent{Erasmus Exchange Master student 1 semester}{Technische Universität München (TUM)}{Oct 2019 - March 2020}{\textbf{1.27} maximum 1 scale 5-1}



% use ONLY \newpage if you want to force a page break for
% ONLY the currentc column
\newpage

\cvsection{Projects}
\begin{itemize}
 \item \textbf{Current Research: Specification Strength}

I am pursuing my Ph.D. under the supervision of Professor Alexandra Mendes (Faculty of Engineering of the University of Porto), co-supervised by Professor Ruben Martins (Carnegie Mellon University).

My current work focuses on automatic methods to assess \textbf{how strong a specification is}, i.e., whether a verified program is actually proving what the developer intended. I lead \textbf{ProofPulse} (ASE 2026), an interactive proof-coverage analysis for Dafny: a three-valued model classifies each specification element as essential, auxiliary or irrelevant to the proof, exposing \textbf{vacuous proofs} and \textbf{unnecessary preconditions} (evaluated on 252 programs, strengthened with unsat-core minimization).

\textbf{Ongoing work, SeVerus:} a context-aware benchmark for \textbf{Verus} proof generation. Starting from verified open-source \textbf{Rust} projects, it extracts the dependency graph between proof artifacts (lemmas, invariants, assertions) and builds sequences of tasks in which agents must write and prove the missing \textbf{auxiliary lemmas} themselves, instead of proving a single isolated lemma. Preliminary results show it produces tasks that current coding agents cannot yet solve. Preprint coming soon.

 \item \textbf{Verified Rust: human-eval-verus.} Contributed \textbf{11 merged pull requests} to \underline{\href{https://github.com/secure-foundations/human-eval-verus}{human-eval-verus}}, the Verus translation of HumanEval: verified Rust implementations with specifications, loop invariants and proofs for 10 problems (40, 46, 77, 96, 110, 120, 121, 126, 133, 151) and a partial proof for problem 47.

 \item \textbf{VeriFixer (FCT project):} researcher on an open-source tool that automatically repairs errors reported by verifiers, starting with Dafny (INESC TEC, INESC-ID, CMU).

 \item \textbf{MIPI Camera Controllers (Synopsys):} verification lead for the CSI-2 camera controllers: constrained-random SystemVerilog/UVM environments, temporal assertions derived from the MIPI specifications, and formal verification where stronger guarantees were needed.

 \item \textbf{Wind Flow:} \textbf{team leader} (10 of 40 students) implementing a long-range wireless physical layer on an SDR for a mesh of wind sensors.
\end{itemize}

%% Switch to the right column. This will now automatically move to the second
%% page if the content is too long.
\switchcolumn


\cvsection{Accomplishments}

\textbf{Merit Awards}
\begin{itemize}
  \item Received in 2016-2017 and 2017-2018 the \textbf{Merit scholarship} awards by academic excellency at FEUP. 

    \item Received the award \textbf{Incentivo} 2018 at FEUP by being one of the \textbf{top 20} 1st year students.
\end{itemize}


%\textbf{Science Awards}
    
%    \begin{itemize}
%        \item \textbf{Bronze medal} in the National \textbf{Mathematics Olympiads} 2011 and  National \textbf{Chemistry Olympiads} 2015
%         \item Represented Portugal in the International\textbf{ Ibero-American Chemistry Olympiads}  2016.
        %\item 1st place in \textbf{Elídio Pinho} science fair 2015
%    \end{itemize}


\cvsection{Programming}
\cvskill{Verus}{5}
\cvskill{Dafny}{5}
\cvskill{Python}{5}
\cvskill{C}{5}
\cvskill{SystemVerilog/UVM}{5}
\cvskill{Lean}{3.5}


\cvsection{Academic Output}

\begin{itemize}
\item
Silva, Álvaro F., Ruben Martins, and Alexandra Mendes. ``ProofPulse: Interactive Proof Coverage Analysis for Dafny.'' \textbf{ASE 2026}.
\item
Pinto, Sofia Vieira, Álvaro F. Silva, João Pascoal Faria, and Alexandra Mendes. ``DSpec2Test: Specification-Driven Test Generation in Dafny.'' \textbf{ASE 2026}.
\item
Silva, Álvaro, Isabel Amaral, João Pascoal Faria, and Alexandra Mendes. ``Improving Debugging in Verification-Aware Languages Through Automated Fault Localization: A Case Study in Dafny.'' \textbf{ISSRE 2026}.
\item
Amaral, Isabel, Álvaro F. Silva, João F. Ferreira, and Alexandra Mendes. ``DafnyFix: Single-Transformation Automated Repair in Dafny.'' \textbf{APR Workshop @ ASE 2026}.
\item
Silva, Álvaro, Alexandra Mendes, and Ruben Martins. ``Inferring Multiple Helper Dafny Assertions with LLMs.'' \textbf{arXiv preprint} 2511.00125 (2025).
\item
Carreira, Carolina, Álvaro Silva, Alexandre Abreu, and Alexandra Mendes. ``Can Large Language Models Help Students Prove Software Correctness? An Experimental Study with Dafny.'' \textbf{SEFM 2025}.
\item
Silva, Álvaro, Alexandra Mendes, and João F. Ferreira. ``Leveraging Large Language Models to Boost Dafny's Developers Productivity.'' \textbf{FormaliSE 2024}.
\end{itemize}

\textbf{Thesis}
\begin{itemize}
    \item Silva, Álvaro. "Using Dimensional Aware Genetic Programming to find interpretable Dispatching Rules for the Job Shop Scheduling Problem". Faculty of Engineering of the University of Porto. (2021).
\end{itemize}


\cvsection{Service}
\begin{itemize}
    \item \textbf{Reviewer}, Symposium on Formal Methods (FM 2026)
    \item \textbf{Artifact Evaluator}, FormaliSE 2026
\end{itemize}


\cvsection{Supervision \& Talks}

\textbf{Thesis Supervision}
\begin{itemize}
\item
de Jesus, Renato Pedrosa. "Evaluating High-Level SystemC Formal Verification Compared to RTL Verification." (2024) (Company supervisor of the thesis).
\end{itemize}

\textbf{Conference Presentations}
\begin{itemize}
\item Silva, Álvaro, Bruno Maurício. "A Simple Approach to Lamport's State Machine Replication". IJUP'2021- 14º edition. Oral communication. (2021)
\end{itemize}


\cvsection{References}

%\cvref{name}{email}{mailing address}
\divider

\cvref{Prof.\ Alexandra Mendes }{INESC TEC and FEUP}{afmendes@fe.up.pt}

\divider

\cvref{Prof.\ Ruben Martins }{Carnegie Mellon}{rubenm@cs.cmu.edu}

\divider

\cvref{Manager ASIC Digital Design.\ Rafael Franco}{Synopsys}{arr.franco@outlook.com}

\divider
\cvref{Prof.\ Gonçalo Reis Figueira}{INESC TEC and FEUP}{lfigueira@fe.up.pt}


\end{paracol}

\end{document}
